\section{Problem 5}
\label{sec:problem5}

The position and inertial velocity of an asteroid relative to the Sun are given in canonical units ($\mu=1$) as
\[
[\vect{r}]_S = \begin{bmatrix} 0.7 \\ 0.6 \\ 0.3 \end{bmatrix}\ \text{AU}, \qquad
[{}^{S}\vect{V}]_S = \begin{bmatrix} -0.8 \\ 0.8 \\ 0 \end{bmatrix}\ \text{AU}\cdot\text{TU}^{-1}.
\]
The six orbital elements (a) are obtained with the Problem 1 function. The period, semi-latus rectum, specific angular momentum, and specific mechanical energy, (b) through (e), then follow from
\begin{equation}
T = 2\pi\sqrt{\frac{a^3}{\mu}}, \qquad
p = a(1-e^2), \qquad
{}^{S}\vect{h} = \vect{r}\times{}^{S}\vect{V}, \quad h = \|{}^{S}\vect{h}\|, \qquad
\En = -\frac{\mu}{2a}.
\end{equation}

\subsection*{Impact check}

Earth is assumed to move on a circular orbit of radius 1 AU in the reference (ecliptic) plane. The asteroid can only cross Earth's orbit at its ascending and descending nodes, where $\omega+\nu = 0^\circ$ and $180^\circ$. The radius at each node follows from the orbit equation,
\begin{equation}
r_{\text{asc}} = \frac{p}{1+e\cos(-\omega)}, \qquad
r_{\text{desc}} = \frac{p}{1+e\cos(180^\circ-\omega)},
\end{equation}
and an impact is possible only if a node lies within one Earth radius of Earth's orbit:
\begin{equation}
\left|r_{\text{node}} - 1\right| \le R_E = \frac{6378}{1.496\times10^{8}}\ \text{AU}.
\end{equation}

\programoutput{outputs/problem5.txt}

\noindent Both nodes lie outside Earth's orbit, at 1.291 and 1.115 AU, so
\begin{equation}
\boxed{\text{the asteroid cannot impact the Earth.}}
\end{equation}

\FloatBarrier
